\documentclass[10pt,a4paper]{article}
\usepackage{amsmath,amssymb,amsthm,graphicx}
\newcommand{\bits}{\{0,1\}}
\newcommand{\DT}{\operatorname{DT}}
\newcommand{\DTSize}{\operatorname{DTSize}}
\newcommand{\CC}{\operatorname C}
\newcommand{\comp}{\diamond}
\newcommand{\OR}{\operatorname{OR}}
\newcommand{\AND}{\operatorname{AND}}
\newcommand{\XOR}{\operatorname{XOR}}
\newenvironment{originalresult}[1]{\par\medskip\noindent\textbf{#1.}\itshape}{\par\medskip}
\title{Complete gadget dichotomy for decision-tree depth-to-size lifting}
\author{Yaroslav Alekseev, Yuval Filmus, Alexander V. Smal}
\date{doi:10.1007/s00037-025-00276-5; CC BY 4.0}
\begin{document}
\maketitle
\noindent\textit{Editorial note.} Faithful editable TeX transcription of the actual CC BY4.0 complete journal version, DOI10.1007/s00037-025-00276-5. CCC abstract is only a source alias and its license is not used to extend to omitted proofs. Complete supporting author random-projection proof of Theorem3.2 is included; its alternative simulation proof is not required and not counted. Complete Theorem3.4, Corollary3.5, Lemma3.8 and Theorem3.9 author proofs retained. Original source notation and wording preserved; standard background is not recursively reproved.\par\medskip
\noindent\textit{Allowed standard prerequisites (editorial).}\begin{itemize}
\item Decision-tree depth versus certificate complexity. DT(f) <= C(f)\textasciicircum{}2. Beals et al.2001, exact statement retained in source Corollary3.5 proof.
\end{itemize}\medskip
\section*{Block composition (author source p.2)}
For $f:\bits^n\to V$ and $g:\bits^m\to\bits$, a (block-)composition $f\comp g:\bits^{n\times m}\to V$ is defined by
\[(f\comp g)(z_1,z_2,\ldots,z_n):=f(g(z_1),g(z_2),\ldots,g(z_n)),\]
where each $z_i\in\bits^m$.
\section*{Prerequisites (author source pp.10--12)}
All the logarithms are base 2. For $g:\bits^k\to\bits$, we define a function $g^n:\bits^{n\times k}\to\bits^n$ such that
\[g^n(x_1,x_2,\ldots,x_n):=(g(x_1),g(x_2),\ldots,g(x_n)).\]
For a function $f:\bits^n\to\bits$, an input $x\in\bits^n$, and a partial assignment $\sigma\in\{0,1,*\}^n$ that agrees with $x$ on all non-$*$ coordinates, we use $f|_\sigma$ to denote a restriction of $f$ to $\sigma$ and $x|_\sigma$ to denote a projection of $x$ to the $*$-coordinates of $\sigma$.
\paragraph{Decision tree complexity.} A decision tree is a rooted binary tree with internal nodes labeled by input variables (representing queries), edges labeled with 0 and 1, and leaves labeled by some values. Decision tree evaluation is defined in a natural way: given an assignment for the input variables, we traverse the tree starting from the root and then sequentially choose every next edge according to the given assignments until we reach some leaf. The label of this leaf is the result of the evaluation. A decision tree computes some function $f$ if for all possible assignments to the input variables, the decision tree evaluates to the value of the function. For a function $f$, $\DT(f)$ denotes the minimal depth of a decision tree computing $f$, and $\DTSize(f)$ denotes the minimal number of leaves in a decision tree computing $f$.
\paragraph{Certificate complexity.} A certificate complexity is a non-deterministic analogue of the decision tree complexity. A certificate for an input $x\in\bits^n$ to a function $f$ is a set $S\subseteq[n]$ of indices such that $f$ restricted to all inputs that match $x$ on $S$ is constant, i.e., $f(y)=f(x)$ whenever $y|_S=x|_S$. The certificate complexity $\CC(f,x)$ of $f$ at input $x$ is the size of the smallest certificate for $x$. Finally, the certificate complexity of function $f$ is defined as $\CC(f):=\max_{x\in\bits^n}\CC(f,x)$.

\section*{3.1 Resistant gadgets (author source pp.14--15)}
\begin{originalresult}{Definition 3.1}
A gadget $g:\bits^m\to\bits$ is resistant if for every $i\in[m]$ and $b\in\bits$, the function obtained by fixing the $i$th input to $b$ is not constant. Equivalently, the minimum certificate complexity at inputs is larger than 1. E.g., the XOR function is resistant while the AND function is not.
\end{originalresult}
\begin{originalresult}{Theorem 3.2}
For any function $f:\bits^n\to\bits$ and a resistant gadget $g:\bits^m\to\bits$,
\[\log\DTSize(f\comp g)=\Omega(\DT(f)).\]
\end{originalresult}
In the following proofs, we use $x_1,\ldots,x_n$ to refer to the inputs of function $f$. The inputs to the lifted function $f\comp g$ are denoted by $z_{i,j}$, where $i\in[n]$ and $j\in[m]$.
\begin{proof}[Proof using random projections]
Let $T$ be a decision tree for $f\comp g$ of size $S$. We apply the following random projection. For each $i\in[n]$, choose an index $j\in[m]$ at random, and set $z_{i,j}$ to a random value $b\in\bits$. Set the remaining variables $z_{i,k}$, $k\ne j$, to $x_i,\bar x_i,0,1$ in such a way that $g(z_{i,1},\ldots,z_{i,m})=x_i$. This is possible since $g$ is resistant, as we show below. The result is a decision tree $T'$ for $f$.
To see that we can always construct such a projection, let $h(z)$ be a non-constant function, and choose two inputs $w^0,w^1$ such that $h(w^b)=b$. We define the projection $\rho$ as follows: if $w_i^0=w_i^1=a$ then set $z_i=a$; otherwise, if $w_i^0=0$ then set $z_i=x$, and if $w_i^1=0$ then set $z_i=\bar x$. By construction, $\rho|_{x=b}=w^b$ and so $h|_\rho(b)=h(w^b)=b$.
For example, suppose that $h(z_1,z_2,z_3)$ is the majority function on three inputs and that we set the first bit $z_1$ to 0. Then $h(0,1,x_i)=x_i$ and $h(0,x_i,x_i)=x_i$ are two possible ways to choose a pair of inputs $(w^0,w^1)$ corresponding to $(010,011)$ and $(000,011)$, respectively.
We would like to say that there is a choice of random projection for which $T'$ is shallow. To do this, let us calculate the probability that some vertex $v$ at depth $D$ in $T$ survives in $T'$. The path leading to $v$ queries variables from at least $D/m$ different blocks. For each such block, choose an arbitrary variable that is queried along the path and consider the value the path assigns to it. If the random projection sets this variable to the opposite value then the vertex $v$ does not survive in $T'$. It happens with probability at least $1/(2m)$ for each variable, and so the vertex $v$ survives with probability $p\le(1-1/(2m))^{D/m}=1/2^{O(D)}$.
If $pS<1$, then the expected number of vertices at depth $D$ surviving in $T'$ is less than 1, and so there is a random projection for which $T'$ has a depth less than $D$. This happens for some $D=O(\log S)$.
\end{proof}

\section*{3.2 Weakly resistant gadgets (author source pp.16--17)}
There are gadgets, such as $x\vee(y\wedge z)$, which are clearly not resistant, but for which the lifting still holds as we will show below. To capture these cases, we define a more general class of weakly resistant gadgets.
\begin{originalresult}{Definition 3.3}
A gadget $g:\bits^m\to\bits$ is weakly resistant if for every certificate $\alpha$ (a partial assignment which sets the value of $g$) there is a substitution of one input variable (with 0 or 1) which conflicts with $\alpha$ and does not make $g$ constant.
\end{originalresult}
Every resistant gadget is trivially weakly resistant: we can take any variable mentioned in $\alpha$ and substitute the opposite value. The aforementioned gadget $h(x,y,z)=x\vee(y\wedge z)$ is weakly resistant (as we show below) but not resistant, since $h|_{x=1}\equiv1$. The gadget $x\vee y$ is not even weakly resistant due to the certificate $x=y=0$: if we substitute $x=1$ or $y=1$ then the gadget becomes constant.
Let us verify that $h$ is weakly resistant by considering all possible certificates:
\begin{itemize}
\item $x=1$: take $x=0$.
\item $y=z=1$: take $y=0$ or $z=0$.
\item $x=y=0$: take $y=1$.
\item $x=z=0$: take $z=1$.
\end{itemize}
We now prove the following lifting theorem:
\begin{originalresult}{Theorem 3.4}
For any $f:\bits^n\to\bits$ and a weakly resistant gadget $g:\bits^m\to\bits$ the following holds:
\[\log\DTSize(f\comp g)=\Omega(\CC(f)).\]
\end{originalresult}
\begin{proof}
Given a decision tree $T$ for $f\comp g$ of size $S$, we construct a blockwise random projection $\rho$ in the following way. For each $i\in[n]$, we choose a random certificate $C_i$ of $g$. Let $z_{i,j}=b$ be the substitution promised by weak resistance for $C_i$ that does not make $g(z_i)$ constant. We can complete $\rho|_{z_i}$ to a projection such that $g(\rho|_{z_i})=x_i$ (as in the case of resistant predicates, see proof of Theorem 3.2 using random projections).
Let $\ell$ be a leaf of $T$ and $\tau$ be the corresponding partial assignment to the variables of $f\comp g$. Suppose that $\tau$ contains certificates for at least $D$ blocks (i.e., in these blocks, the value of $g$ on $\tau$ is fixed). If there are $M$ certificates of $g$, then $\tau$ is consistent with $\rho$ with probability at most $(1-1/M)^D$. Choose $D=\Theta(\log S)$ so that $(1-1/M)^DS<1$. By union bound, under some choice of $\rho$, the partial assignments for all leaves of $T|_\rho$ contain certificates for less than $D$ blocks. Since by construction $(f\comp g)|_\rho=f$, this shows that $\CC(f)<D$.
Indeed, given an input $x$ to $f$, consider the corresponding input $\rho(x)$ to $f\comp g$. This input reaches some leaf $\ell$, in which less than $D$ blocks contain certificates, say the blocks in $I\subseteq[n]$. For each input $x'$ consistent with gadget values for blocks in $I$, we can find an input $z$ reaching the same leaf such that $g^n(z)=x'$. Hence, the block certificates in $\ell$ correspond to a certificate of $f$.
\end{proof}
\paragraph{Remark.} One of the anonymous reviewers noticed that the structure of a decision tree is completely irrelevant to the above proof. We only use the fact that the leaves are monochromatic and cover all the inputs. Thus, Theorem 3.4 can be restated by replacing the decision tree size with the subcube cover number. This also explains why, in the next section, we will be able to reuse the above proof in a similar lifting result for conjunction DAGs (see Theorem 4.2).
\begin{originalresult}{Corollary 3.5}
For any $f:\bits^n\to\bits$ and a weakly resistant gadget $g:\bits^m\to\bits$ the following holds:
\[\log\DTSize(f\comp g)=\Omega(\sqrt{\DT(f)}).\]
\end{originalresult}
\begin{proof}
Beals et al. (2001) shows that for any function $f$,
\[\DT(f)\le(\CC(f))^2.\]
Together with Theorem 3.4, that gives us the desired bound.
\end{proof}

\section*{3.3 Gadget classification (author source pp.18--19)}
Now, we can prove a dichotomy result. First of all, we need the following classification of gadgets.
\begin{originalresult}{Lemma 3.8}
For every $g:\bits^m\to\bits$, one of the following cases holds:
\begin{enumerate}
\item[(i)] Function $g$ is a (possibly empty) conjunction or a disjunction of literals.
\item[(ii)] Function $g$ is weakly resistant.
\end{enumerate}
\end{originalresult}
\begin{proof}
Let $y_1,\ldots,y_m$ denote the inputs of $g$. If $g$ is not weakly resistant then there is a certificate $\alpha$ such that for every substitution $y_j=b$ conflicting with $\alpha$, the function $g|_{y_j=b}$ is constant. Moreover, we can assume that $\alpha$ is an inclusion-minimal certificate.
Suppose that $\alpha=\{y_{i_1}=b_1,\ldots,y_{i_t}=b_t\}$ and $g|_\alpha=b$. By assumption, the function $g|_{y_{i_s}=\bar b_s}$ is constant. Since $\alpha$ is inclusion-minimal, necessarily $g|_{y_{i_s}=\bar b_s}=\bar b$. Thus if $b=1$ then
\[g=[y_{i_1}=b_1]\wedge\cdots\wedge[y_{i_t}=b_t],\]
otherwise, if $b=0$ then
\[g=[y_{i_1}=b_1]\vee\cdots\vee[y_{i_t}=b_t].\]
\end{proof}
We can now state the dichotomy theorem.

\begin{originalresult}{Theorem 3.9}
For every gadget $g:\bits^m\to\bits$, one of the following cases holds:
\begin{enumerate}
\item[(i)] There is an infinite family of functions $f_n$ with $\DT(f_n)\to\infty$ and $\DTSize(f_n\comp g)=O(\DT(f_n))$.
\item[(ii)] For every function $f$, $\log\DTSize(f\comp g)=\Omega(\DT(f)^{\Omega(1)})$.
\end{enumerate}
\end{originalresult}

\begin{proof}
These two cases correspond to the cases of Lemma 3.8. If $g$ is a disjunction then we take $f_n:=\OR_n$. Clearly $\DT(f_n)=n$ and $\DTSize(f\comp g)=O(n)$ (if $g$ is constant, then $\DTSize(f_n)=1$). The case of a conjunction is similar.
Otherwise, $g$ is weakly resistant, thus by Corollary 3.5 we have $\log\DTSize(f\comp g)=\Omega(\sqrt{\DT(f)})$.
\end{proof}

\begin{thebibliography}{99}
\section*{Source reference used as standard background}
Robert Beals, Harry Buhrman, Richard Cleve, Michele Mosca and Ronald de Wolf (2001). Quantum Lower Bounds by Polynomials. J. ACM 48(4), 778--797. doi:10.1145/502090.502097.
\end{thebibliography}
\end{document}
